\documentclass[11pt,a4paper]{article}

\usepackage[T1]{fontenc}
\usepackage[utf8]{inputenc}
\usepackage[brazil]{babel}
\usepackage{lmodern}
\usepackage{microtype}
\usepackage[margin=2cm]{geometry}
\usepackage{hyperref}
\usepackage{titlesec}
\usepackage{xcolor}

\hypersetup{
  colorlinks=true,
  linkcolor=black,
  urlcolor=black,
  citecolor=black,
  pdfborder={0 0 0}
}

\definecolor{sectioncolor}{HTML}{111111}
\definecolor{accentcolor}{HTML}{1A1A1A}

\pagestyle{empty}
\setlength{\parindent}{0pt}
\setlength{\parskip}{5pt}
\linespread{1.05}
\selectfont

\titleformat{\section}
  {\bfseries\large\color{sectioncolor}}
  {}{0pt}{}[\vspace{-2pt}{\color{sectioncolor}\titlerule[0.8pt]}]
\titlespacing*{\section}{0pt}{12pt}{8pt}

\newcommand{\cvrole}[4]{%
  \vspace{4pt}%
  {\bfseries\color{accentcolor}#1}\hfill{\itshape #2}\\[-2pt]%
  {\itshape #3}\hfill #4\\[6pt]%
}

\begin{document}

\begin{center}
  {\fontsize{20}{24}\selectfont\bfseries Kamille Hartmann Maccagnan}\\[8pt]
  {\large Analista de BI Júnior \textbar{} Power BI, SQL \& Excel \textbar{} KPIs}\\[10pt]
  Campo Comprido, Curitiba-PR \enspace\textbullet\enspace (41) 9 9928-4424 \enspace\textbullet\enspace
  \href{mailto:kamillehartmannm@gmail.com}{kamillehartmannm@gmail.com}\\[3pt]
  \href{https://linkedin.com/in/kamille-hartmann-maccagnan/}{linkedin.com/in/kamille-hartmann-maccagnan/}
\end{center}

\vspace{6pt}

\section{Resumo Profissional}

Analista de Business Intelligence com experiência na construção e manutenção de dashboards em Power BI, desenvolvimento de medidas em DAX e consolidação de informações para apoio à tomada de decisão em ambientes operacionais, corporativos e industriais. Atua em consultas e extração de dados com SQL, preparação e transformação com Power Query e Excel, monitoramento de KPIs, padronização de informações e governança da qualidade dos dados. Vivência em interface com áreas de negócio, operações e equipes técnicas, elaboração de relatórios gerenciais e apresentações executivas. Perfil analítico, organizado, proativo e colaborativo, com inglês avançado.

\section{Experiência Profissional}

\cvrole{Anjuss}{Analista de Suporte de TI}{07/2025 -- Atual}

Desenvolvo e mantenho dashboards e relatórios em Power BI para monitoramento de indicadores operacionais e de processos, apoiando a consolidação de informações gerenciais e a tomada de decisão baseada em dados. Realizo extração, preparação e transformação de dados com Power Query e Excel, com validação para garantir integridade, consistência e confiabilidade das informações. Elaboro e atualizo documentações e POPs, apoiando a padronização de fluxos e procedimentos. Atuo como ponte entre usuários de negócio e áreas técnicas, alinhando demandas, acompanhando entregas e identificando oportunidades de melhoria e automação nos relatórios existentes.

\cvrole{HORSE do Brasil}{Estagiária de TI (IS/IT LATAM)}{07/2024 -- 07/2025}

Atuei em ambiente industrial multinacional apoiando equipes de dados e PMO na elaboração de relatórios executivos e dashboards em Power BI para acompanhamento de performance de projetos, operações e iniciativas estratégicas. Participei de processos de governança de dados, documentação, verificação de qualidade das informações e apoio a ETL com Power Query para análises confiáveis. Implementei fluxos no Power Automate para distribuição automatizada de relatórios de BI. Colaborei com áreas de negócio, operações e stakeholders internacionais em reuniões multilíngues, facilitando comunicação, alinhamento de demandas e entrega de materiais de apoio à gestão.

\cvrole{Restaurante Familiar}{Analista de Dados e Suporte Técnico}{01/2023 -- 07/2024}

Estruturei do zero a base de dados operacional em Excel e desenvolvi dashboards em Power BI para consolidação e monitoramento de vendas, resultados financeiros e indicadores de desempenho. Realizei consultas, integração e padronização de dados de múltiplas fontes operacionais, com validação crítica para garantir qualidade e precisão das análises. Modelei KPIs e organizei controles e relatórios gerenciais para apoio à gestão, identificação de desvios e oportunidades de melhoria contínua nos processos operacionais e administrativos.

\section{Competências Técnicas}

\textbf{Business Intelligence:} Power BI (DAX, Power Query), desenvolvimento e manutenção de dashboards, relatórios gerenciais e executivos, visualização de dados, KPIs, modelagem de dados, storytelling com dados.

\textbf{Dados e ETL:} SQL, Excel avançado, Power Query, ETL/ELT, bancos de dados relacionais, consolidação de informações, validação e governança de dados, qualidade e precisão das informações.

\textbf{Processos e Comunicação:} Padronização de procedimentos, documentação e POPs, apresentações em PowerPoint, interface com stakeholders, apoio a reuniões e treinamentos, acompanhamento de prazos e entregas, melhoria contínua.

\textbf{Ferramentas:} Power Automate, ServiceNow, Google Sheets, Trello, Monday, Pacote Office.

\textbf{Metodologias:} Scrum, Kanban, Boas Práticas PMOBOK.

\textbf{Idiomas:} Inglês Avançado \enspace\textbullet\enspace Espanhol Básico \enspace\textbullet\enspace Coreano Básico

\section{Projetos e Conquistas}

1º lugar no Transformation Day Renault 2023 com solução orientada a dados e melhoria de processos. Participação no Hackathon Bosch 2023 com foco em resolução de problemas de negócio com dados. Implementação ponta a ponta de bases em Excel e dashboards em Power BI para monitoramento de performance operacional. Automação de fluxos de distribuição de relatórios de BI com Power Automate integrado ao Power BI na HORSE.

\section{Formação Acadêmica}

\textbf{PUCPR} \hfill Curitiba, PR\\
Graduação em Gestão da Tecnologia da Informação \hfill Previsão de conclusão: 06/2026\\[6pt]

\textbf{Escola Conquer} \hfill Curitiba, PR\\
Pós-graduação em Liderança e Gestão em Tecnologia \hfill Conclusão: 2024\\[6pt]

\textbf{Universidade Tuiuti do Paraná} \hfill Curitiba, PR\\
Graduação em Medicina Veterinária \hfill 06/2019\\[6pt]

\textbf{Qualittas} \hfill Curitiba, PR\\
Pós-graduação em Clínica Médica e Cirúrgica de Animais Exóticos e Pets Não Convencionais \hfill 12/2022

\end{document}
